% pIII_canonical.tex --- Phase B deliverable
% Session VQUAD-PIII-NORMALIZATION-MAP, 2026-05-02

\section*{Canonical $P_{III}(D_6)$ Hamiltonian form}
\label{sec:pIII-canonical}

\paragraph{Convention.}
We adopt the Okamoto $\sigma$-form Hamiltonian for
$P_{III}(D_6)$ (Okamoto 1987, ``Studies on the Painlev\'e
equations.~IV. Third Painlev\'e equation $P_{III}$'',
\textsl{Funkcial.~Ekvac.}~30, 305--332; reproduced in
Conte--Musette 2008, Ch.~7), which has parameters
$(\alpha_\infty,\alpha_0,\beta_\infty,\beta_0)$ subject to the
constraint
\[
  \alpha_\infty + \alpha_0 + \beta_\infty + \beta_0 = 0.
\]
The Hamiltonian, as quoted in the relay-prompt 009 brief
(consistent with Okamoto--Sakai normalization), reads
\[
  \boxed{\;
    t\,H_{III(D_6)}(q,p;t) \;=\;
      2\,q^2 p^2
      \;-\;\bigl(2 t\,q^2 + (2\alpha_\infty+1)\,q - 2 t\bigr)\,p
      \;+\;(\alpha_\infty + \alpha_0)\,t\,q,
  \;}
\]
with Hamilton's equations $t\,q' = \partial(t H)/\partial p$,
$t\,p' = -\partial(t H)/\partial q$:
\[
  \begin{aligned}
    t\,q' &= 4\,q^2 p - 2\,t\,q^2 - (2\alpha_\infty+1)\,q + 2\,t, \\
    t\,p' &= -4\,q\,p^2 + 4\,t\,q\,p + (2\alpha_\infty+1)\,p
              - (\alpha_\infty+\alpha_0)\,t.
  \end{aligned}
\]
Eliminating $p$ yields the standard $P_{III}$ equation in $q(t)$
\[
  q'' = \tfrac{(q')^2}{q} - \tfrac{q'}{t}
       + \tfrac{1}{t}\bigl(\alpha\,q^2 + \beta\bigr)
       + \gamma\,q^3 + \tfrac{\delta}{q},
\]
with the standard $(\alpha,\beta,\gamma,\delta)$ parameters
related to $(\alpha_\infty,\alpha_0,\beta_\infty,\beta_0)$ by
the Okamoto--Sakai dictionary.

\paragraph{Singular structure of $P_{III}(D_6)$.}
The equation has irregular singularities at $t=0$ and
$t=\infty$, both of Poincar\'e rank $1$ (level-$1$ trans-series
in $1/t^{1/2}$ at $t=0$ and in $t^{1/2}$ at $t=\infty$). The
$D_6$ refinement specifies the local-monodromy data at $t=0$
to be of toric type with full irregular structure (as opposed
to $P_{III}(D_7)$ and $P_{III}(D_8)$ degenerations).

\paragraph{V$_\mathrm{quad}$ parameter point.}
CT~v1.3 §3.5, \textsf{Thm.~vquad-cc}, identifies the
V$_\mathrm{quad}$ family with $P_{III}(D_6)$ at the parameter
point
\[
  (1/6,\ 0,\ 0,\ -1/2)
  \quad\text{in some convention of }
  (\alpha_\infty,\alpha_0,\beta_\infty,\beta_0).
\]
\textsc{Caveat.} The mapping
$(\text{CT v1.3 4-tuple})\to(\alpha_\infty,\alpha_0,\beta_\infty,\beta_0)$
is \emph{not} pinned by any source local to the operator's
library: CT~v1.3 \textsf{Thm.~vquad-cc} cites the parameter
point but defers the explicit Okamoto-convention
identification to \textsf{op:cc-formal-borel}. This is exactly
the residual G15 question.

\paragraph{Constraint check.}
The numerical values $\{1/6, 0, 0, -1/2\}$ sum to
$1/6 - 1/2 = -1/3 \neq 0$. They therefore do \emph{not}
directly satisfy the Okamoto-convention constraint
$\alpha_\infty+\alpha_0+\beta_\infty+\beta_0=0$. This means
either (a) CT~v1.3 §3.5 uses a different convention (e.g.\
the Sakai $E_7^{(1)}/D_7^{(1)}$ root-data convention, where
the four parameters are not constrained to sum to zero), or
(b) one of the four entries should be reinterpreted as a
spectral-type label rather than a Hamiltonian parameter.

\textbf{This convention question is residual~R1 of G15.}
Without Okamoto~1987 §2 or Conte--Musette ch.~7 §7.3 we cannot
disambiguate.

\paragraph{Tau function and Stokes data.}
The $P_{III}(D_6)$ tau function $\tau(t)$, with
$H_{III(D_6)} = (\log\tau)'$, has connection (Stokes) data
between the $t=0$ irregular sector and the $t=\infty$
irregular sector encoded by a pair of Stokes multipliers
$(s_0, s_\infty)$ taking values on a complex algebraic surface
(the cubic surface of monodromy data; see Lisovyy--Roussillon
2017 for a modern treatment). The literature's standard
``canonical'' Stokes-data normalization fixes:
\begin{enumerate}
  \item the multiplicative scale of the Lax-pair connection
        matrix at $t=0$ (a free factor in the formal
        fundamental solution),
  \item the relative phase between the $t=0$ and $t=\infty$
        sectors (a $\mathrm{U}(1)$ ambiguity that is fixed by
        Boutroux's principal asymptotic value),
  \item the labeling of the two Stokes rays at each irregular
        sector (combinatorial $\mathbb{Z}/2$).
\end{enumerate}
The map from \emph{V$_\mathrm{quad}$-native} Stokes data
$(C, \zeta_*)$ to canonical $P_{III}(D_6)$ Stokes data
$(s_0^{\mathrm{can}}, s_\infty^{\mathrm{can}})$ is precisely
$\Phi$ of \texttt{phi\_change\_of\_variables.tex}.
